El código Mermaid devuelto inicialmente por la IA fue sintácticamente correcto y compilable, respetando los cuatro participantes (\texttt{U}, \texttt{FE}, \texttt{BE}, \texttt{DB}) y el uso de flechas síncronas (\texttt{->>}) y de respuesta (\texttt{-->>}). Sin embargo, la auditoría técnica reveló tres puntos de mejora que se corrigieron manualmente:
\begin{itemize}[leftmargin=*,nosep]
\item \textbf{Endpoints y parámetros en inglés ajenos al glosario común:} La IA utilizó \texttt{GET /api/products/:id} y \texttt{POST /api/cart/items (productId, size, quantity)}. Se reemplazaron por las rutas normalizadas del proyecto: \texttt{GET /api/productos/id} y \texttt{POST /api/carrito/lineas (idProducto, talla, cantidad)}.
\item \textbf{Sustitución de frases descriptivas por métodos del Diagrama de Clases:} En el Backend se reemplazaron los textos en prosa de la IA por invocaciones directas a los métodos definidos en \texttt{clases.mmd}: \texttt{findProductoById(id)}, \texttt{findStockByProducto(id)}, \texttt{Producto.tallasDisponibles()}, \texttt{StockTalla.hayStock(cantidad)}, \texttt{Carrito.agregarProducto(p, talla, cantidad)}, \texttt{LineaCarrito.calcularImporte()} y \texttt{Carrito.calcularTotal()}.
\item \textbf{Eliminación de anidamiento \texttt{404} superfluo y ajuste de cabecera:} Se eliminó el bloque exterior \texttt{alt Producto encontrado / else 404 Not Found} (ya que el cliente accede desde una tarjeta existente del catálogo), centrando el bloque condicional en \texttt{alt hayStock(cantidad) == true / else hayStock(cantidad) == false (409 Conflict)}, y se incorporó \texttt{mirrorActors: false}.
\end{itemize}
